\section{Worum es geht}

\textbf{[SETZUNG]} Dieses Dokument schließt die Lücke zwischen der
Halbzähligkeit auf dem Torus (A263) und der chiralen Projektion in den
flachen Raum (A090, A264). Der Kern: ein Torus hat eine Innen- und eine
Außenseite mit topologisch verschiedenen Kopplungsrichtungen. Nach der
Dekompaktifizierung erscheinen diese als zwei chirale Zustände -- und
die Projektion ist nicht symmetrisch.

\section{Halbzähligkeit und Torusgeometrie}

\textbf{[B]} Spin-$\tfrac12$ folgt aus der Wicklung $(n_\theta,n_\phi)=(1,1)$
auf $T^4$ (A263): zwei Umdrehungen $(720°)$ bringen den Spinor in den
Ausgangszustand. Das ist reine Topologie.

\textbf{[B]} Ein Torus mit Wicklung $(1,1)$ hat eine intrinsische
Auszeichnung: Innen- und Außenseite. Für einen eingebetteten Torus im
$\mathbb{R}^3$ gilt: die Gaußsche Krümmung ist außen positiv ($K>0$) und
innen negativ ($K<0$). Für den flachen Identifikations-Torus $T^4$ gilt
$K=0$ überall -- aber die poloidale/toroidale Zirkulationsstruktur definiert
trotzdem eine topologische Innen/Außen-Unterscheidung durch die Orientierung
der Wicklung.

\textbf{[B]} Die Feldrichtung ist durch die Wicklungsgeometrie topologisch
festgelegt: der elektrische Feldvektor zeigt auf der Außenseite nach innen
(für negative Ladung) bzw. nach außen (für positive Ladung). Diese Richtung
ist nicht frei wählbar -- sie folgt aus der quantisierten Wicklungszahl
$\Phi=n\cdot h/e$.

\section{Chirale Projektion}

\textbf{[B]} Die Dekompaktifizierung $T^4\to\mathbb{R}^3\times\mathbb{R}_t$
(Typ-I-Projektion, A090) ist nicht symmetrisch bezüglich Innen und Außen.
Der kompakte Torus trägt beide Seiten gleichzeitig; der dekompaktifizierte
flache Raum sieht nur eine Projektion.

\textbf{[B]} Was nach der Projektion übrigbleibt:
\begin{itemize}[leftmargin=2em,itemsep=2pt,topsep=2pt]
  \item Der \emph{Außenzustand} projiziert auf linkshändige Kopplung
    (positive Rotationsphase $+\xi/2$).
  \item Der \emph{Innenzustand} projiziert auf rechtshändige Kopplung
    (negative Rotationsphase $-\xi/2$).
\end{itemize}
Die Projektion wählt eine Seite aus -- das ist der geometrische Ursprung
der chiralen Asymmetrie.

\textbf{[B]} \textbf{Warum das die chirale Struktur erklärt.} Im kompakten
$T^4$ koppeln beide Seiten mit gleicher Stärke -- es gibt keine bevorzugte
Richtung. Nach der Dekompaktifizierung existiert nur noch die Außenperspektive
als beobachtbarer Zustand. Der Innenzustand ist durch die Typ-I-Projektion
nicht direkt zugänglich -- er ist der ``andere Zweig'', der bei der
Dimensionsreduktion verloren geht (A090, Typ-II-Verlust).

\section{Verbindung zur schwachen Wechselwirkung}

\textbf{[B]} \textbf{Fourier-Projektor und $g_R=0$.} Der Hilbertraum auf
$T^4$ zerfällt in Sektoren nach dem Wicklungsindex $m$. Der physikalische
Projektor $P_+$ wählt den $m\ge0$-Sektor (Außenzustand):
\[
  P_+ = \sum_{n,\,m\ge0} |n,m\rangle\langle n,m|.
\]
Wirkung auf die beiden chiralen Zustände:
\[
  P_+|1,{+}1\rangle = |1,{+}1\rangle \quad\text{(Außenzustand überlebt)},
  \qquad
  P_+|1,{-}1\rangle = 0 \quad\text{(Innenzustand liegt im Kern)}.
\]
Nach der Projektion ist der Innenzustand $|0\rangle$ -- ein Nullvektor.
Die schwache Kopplung:
\[
  g_L = \langle 1,{+}1|H_\text{weak}|1,{+}1\rangle \neq 0,
  \qquad
  g_R = \langle 0|H_\text{weak}|0\rangle = 0.
\]
\textbf{Das ist kein Vorzeichen-Argument ($g_R=-g_L$), sondern ein
Kern-Argument: der Innenzustand liegt im Kern von $P_+$ und hat nach der
Projektion keine physikalische Existenz mehr.} $g_R=0$ folgt algebraisch.

\textbf{[S]} \textbf{Was noch offen ist: warum $P_+$ und nicht $P_-$?}
Der Projektor $P_+$ entspricht der Außenseite ($K>0$, positive Orientierung,
elektrischer Fluss nach innen). $P_-$ wäre der Innenzustand-Projektor.
Dass $P_+$ der physikalische Projektor ist, folgt aus Orientierungsargumenten
(Gaußsche Krümmung, Feldrotation $\nabla\times\mathbf{E}_\text{field}$) --
aber nicht aus einer ableitbaren Symmetriebrechung. Das ist die eine
verbleibende offene Kante.

\section{Zeichenerklärung}

Symbole die in der gesamten A-Serie verwendet werden, sind in A010 erklärt.
Spezifisch für dieses Dokument:

\begin{center}
\renewcommand{\arraystretch}{1.3}
{\small\begin{tabular}{lp{9cm}}
\toprule
Symbol & Bedeutung \\
\midrule
$(n_\theta,n_\phi)=(1,1)$ & Wicklungszahlen für Spin-$\tfrac12$ auf $T^4$ \\
$K$ & Gaußsche Krümmung; außen $K>0$, innen $K<0$ \\
$\Phi=n\cdot h/e$ & Quantisierte Flusslinien, topol.\ festgelegte Feldrichtung \\
$g_L, g_R$ & Links- und rechtshändige schwache Kopplung \\
Typ-I-Projektion & Informationserhaltende Dekompaktifizierung (A090) \\
Typ-II-Verlust & Nicht-umkehrbare räumliche Projektion, Informationsverlust \\
\bottomrule
\end{tabular}}
\renewcommand{\arraystretch}{1.0}
\end{center}

\section{Prüfskript}

\begin{itemize}[leftmargin=2em,itemsep=2pt,topsep=2pt]
  \item Topologische Konsistenz: Wicklung $(1,1)$, 720°-Symmetrie,
    Gaußsche Krümmung innen/außen, Phasendifferenz $\pm\xi/2$:
    \texttt{python/A\_Serie\_Skripte/a268\_torus\_chiralitaet.py}
\end{itemize}

\section{Was offen bleibt}

\textbf{$g_R=0$ ist jetzt auf [B]-Niveau:} gegeben $P_+$ als physikalischen
Projektor folgt $g_R=0$ algebraisch aus $P_+|1,{-}1\rangle=0$.

\textbf{Was auf [S] bleibt: warum $P_+$ der physikalische Projektor ist.}
Das Orientierungsargument (Gaußsche Krümmung, Feldrotation) ist geometrisch
plausibel, aber nicht aus einer tieferen Symmetriebrechung abgeleitet.
Das ist die verbleibende offene Kante -- und sie ist präzise benannt.

\textbf{Verbindung zum Eichsektor (A190) fehlt.} Die chirale Projektion
erklärt die $SU(2)_L$-Kopplung qualitativ; ob sie die volle
elektroschwache Eichstruktur reproduziert, ist nicht gezeigt.

\section{Ersetzt}

Dieses Dokument systematisiert den Innen/Außen-Chiralitätsmechanismus
aus dem Altkorpus (Teilchenmassen-Torus-Modell, Dirac-Wicklungsstruktur,
Projektionsketten-Framework). Es ergänzt A263 (Dirac, Halbzähligkeit),
A090 (Projektionskette, Typ-I/II) und A264 (Asymmetrie, CP) um das
fehlende Bindeglied.

\section{Verwendung von KI-Werkzeugen}

Dieses Dokument ist unter Verwendung KI-gestützter Werkzeuge entstanden.
Verantwortung für Inhalt und Fehler liegt beim Autor. Wortlaut der Regeln in A010.
